\graphicspath{{Parti/P01-Preliminari/A03-Applicazioni_Lineari_Dualita/Figure/}}
%	
\chapter{Applicazioni lineari e dualità}\label{app:applicazionilin_dualita}



\begin{sommario}
	
	Si introducono i concetti di applicazione fra insiemi, iniettività, suriettività e biettività, inquadrando il \textit{problema inverso} come filo conduttore. Nel caso delle applicazioni lineari fra spazi di dimensione finita, nucleo, immagine e operatore aggiunto determinano i quattro sottospazi fondamentali; le relazioni di ortogonalità fra questi forniscono la condizione di esistenza per il problema inverso e stabiliscono l'uguaglianza fra il rango di un'applicazione e quello del suo aggiunto. Si introduce quindi, accanto al problema primale $\Abb\xb=\bb$, il problema duale $\Abb^\top\yb=\cb$, discutendone la relazione di reciprocità; la classificazione dei due problemi in base alla dimensione dei rispettivi nuclei chiude il capitolo.
\end{sommario}





	%-------------------------------------------------------------------------------------------------------------------
	\section{Applicazioni fra insiemi}
	%-------------------------------------------------------------------------------------------------------------------
	
	Nel capitolo precedente abbiamo studiato i vettori come oggetti a sé stanti, insieme alle operazioni --- somma, prodotto scalare, prodotto vettoriale --- che permettono di combinarli e confrontarli. Un problema meccanico (ad esempio strutturale), però, raramente riguarda un vettore isolato: riguarda piuttosto il \textit{legame} fra grandezze diverse. Si pensi al monitoraggio di un ponte: i sensori installati sull'impalcato restituiscono gli spostamenti effettivamente misurati, e da questi si vuole risalire ai carichi --- traffico, vento, un cedimento imprevisto --- che li hanno prodotti. Oppure si pensi al progetto di un elemento strutturale: si parte da un requisito sulla risposta (uno spostamento massimo ammissibile, una tensione da non superare) e si cerca la configurazione, o l'azione, che lo realizzi. In entrambi i casi si conosce l'effetto e si cerca la causa che lo ha prodotto: si tratta, in generale, di un \textit{problema inverso}.
	
	In termini astratti, se $\Abb$ \`e l'applicazione che a ogni causa $\xb$ associa la risposta $\Abb(\xb)$ del sistema, il problema inverso si scrive: dato $\bb$, trovare $\xb$ tale che $\Abb(\xb)=\bb$. A differenza del problema diretto, che consiste semplicemente nel calcolare $\Abb(\xb)$ per un $\xb$ noto, il problema inverso solleva questioni che nell'altro verso non si pongono nemmeno: una tale $\xb$ esiste sempre? Se esiste, \`e l'unica? La risposta, in generale, \`e negativa su entrambi i fronti --- carichi diversi possono produrre la stessa risposta misurata, cos\`i come un requisito di progetto pu\`o risultare irrealizzabile --- e dipende esclusivamente dalla struttura dell'applicazione $\Abb$, non dal fatto che $\xb$ e $\yb$ rappresentino grandezze meccaniche. \`E dunque naturale isolare queste questioni --- esistenza e unicit\`a della soluzione --- nel contesto pi\`u generale delle applicazioni fra insiemi, prima ancora di specializzarsi al caso lineare: solo dopo aver capito cosa le determina in generale ha senso chiedersi come la linearit\`a le renda trattabili, e solo pi\`u avanti, fissando una base, tradurre tutto in termini di matrici e sistemi di equazioni.

\subsection{Definizioni fondamentali}

Siano $X$ e $Y$ due insiemi non vuoti. Un'\textit{applicazione}\index{applicazione} (o \textit{funzione}) $f: X \to Y$ \`e una
legge che associa a ogni elemento $x \in X$ uno e un solo elemento 	$y = f(x) \in Y$ (ma elementi distinti di $X$ possono avere la stessa
immagine in $Y$).
L'insieme $X$ si chiama \textit{dominio}, l'insieme $Y$ si chiama \textit{codominio}. L'\textit{immagine} di $f$ \`e il sottoinsieme di $Y$ effettivamente raggiunto dall'applicazione:
\begin{equation}
	f(X) = \{y \in Y \mid y = f(x) \text{ per qualche } x \in X\} \subseteq Y.
\end{equation}
In generale, l'immagine $f(X)$ pu\`o essere un sottoinsieme proprio del codominio $Y$.


\subsection{Iniettivit\`a, suriettivit\`a, biettivit\`a}

Un'applicazione $f: X \to Y$ pu\`o godere delle seguenti propriet\`a:
\begin{enumerate}[label=(\roman*)]
	\item $f$ \`e \textit{iniettiva}\index{applicazione!iniettiva} se elementi distinti del dominio hanno immagini distinte:
	\begin{equation}
		x_1 \neq x_2 \;\Longrightarrow\; f(x_1) \neq f(x_2),
	\end{equation}
	ovvero, equivalentemente, $f(x_1)=f(x_2) \Rightarrow x_1=x_2$;
	\item $f$ \`e \textit{suriettiva}\index{applicazione!suriettiva} se l'immagine coincide con il codominio: $f(X) = Y$, ossia per ogni $y \in Y$ esiste almeno un $x \in X$ tale che $f(x)=y$;
	\item $f$ \`e \textit{biettiva}\index{applicazione!biiettiva} (o \textit{biunivoca}) se \`e sia iniettiva sia suriettiva.
\end{enumerate}

\begin{esempio}
	Si consideri il caso $X = Y = \mathbb{R}$.
	L'applicazione $f(x) = 2x+1$ \`e biettiva: \`e iniettiva perch\'e $f(x_1)=f(x_2)$ implica $x_1=x_2$, ed \`e suriettiva perch\'e per ogni $\bar{y} \in \mathbb{R}$ esiste $x=(\bar{y}-1)/2$.
	L'applicazione $f(x) = x^2$ non \`e iniettiva (ad esempio $f(-1)=f(1)=1$) e non \`e suriettiva su $\mathbb{R}$ (i valori negativi non sono raggiunti).
	L'applicazione $f(x) = e^x$ \`e iniettiva ma non suriettiva su $\mathbb{R}$ (l'immagine \`e il solo semiasse positivo).
	L'applicazione $f(x) = x^3-x$ \`e suriettiva ma non iniettiva: \`e suriettiva perch\'e $f(x)\to\pm\infty$ per $x\to\pm\infty$ e, per il teorema dei valori intermedi, ogni valore reale \`e raggiunto; non \`e iniettiva perch\'e, ad esempio, $f(-1)=f(0)=f(1)=0$. I quattro esempi esauriscono le combinazioni possibili di iniettivit\`a e suriettivit\`a, mostrando che si tratta di propriet\`a indipendenti l'una dall'altra.
\end{esempio}


\subsection{Il problema inverso}\label{par:problema_inverso}

Data un'applicazione $f: X \to Y$ e un elemento $\bar{y} \in Y$, il \textit{problema inverso}\index{problema inverso|textbf} consiste nel determinare tutti gli elementi $x \in X$ tali che $f(x) = \bar{y}$. L'insieme di tali elementi \`e la \textit{controimmagine}\index{controimmagine} (o \textit{fibra}) di $\bar{y}$:
\begin{equation}
	f^{-1}(\bar{y}) = \{x \in X \mid f(x) = \bar{y}\}.
\end{equation}
Il problema inverso pone due questioni fondamentali:
\begin{enumerate}[label=(\roman*)]
	\item \textit{Esistenza}: la controimmagine \`e non vuota? Questo \`e garantito se e solo se $\bar{y} \in f(X)$; se $f$ \`e suriettiva, l'esistenza \`e assicurata per ogni $\bar{y} \in Y$.
	\item \textit{Unicit\`a}: la controimmagine contiene al pi\`u un elemento? Per il singolo $\bar{y}$ considerato ci\`o pu\`o valere anche se $f$ non \`e iniettiva altrove --- ad esempio $f(x)=x^2$ ha controimmagine $\{0\}$ per $\bar{y}=0$, pur non essendo iniettiva su $\mathbb{R}$. Ci\`o che garantisce l'unicit\`a per \textit{ogni} $\bar{y}\in Y$ \`e invece l'iniettivit\`a di $f$: $f$ \`e iniettiva se e solo se $f^{-1}(\bar{y})$ contiene al pi\`u un elemento per ogni $\bar{y}\in Y$.
\end{enumerate}
Se $f$ \`e biettiva, il problema inverso ammette un'unica soluzione per ogni $\bar{y} \in Y$ e l'applicazione inversa $f^{-1}: Y \to X$ \`e essa stessa un'applicazione (biettiva).

%\begin{oss}
%	Le questioni di esistenza e unicit\`a del problema inverso coincidono con quelle poste nel \CAPref{app:algebra_eqlineari} per i sistemi di equazioni lineari. L'obiettivo delle prossime sezioni \`e mostrare come, nel caso lineare, la risposta a entrambe le questioni sia completamente determinata dalla struttura dell'operatore.
%\end{oss}



%-------------------------------------------------------------------------------------------------------------------
\section{Applicazioni lineari}\label{sec:applicazioni_lineari}
%-------------------------------------------------------------------------------------------------------------------
La linearità è la proprietà che rende il problema inverso 	trattabile in modo sistematico. Come si vedrà, essa trasforma 	le domande di esistenza e unicità della soluzione --- poste in termini generali nella sezione precedente --- in domande 	sulla struttura di due sottospazi associati all'\textit{operatore} --- termine che, d'ora in poi, useremo come sinonimo di ``applicazione lineare'': il nucleo e l'immagine. Per il momento tratteremo $\Abb$ in modo astratto, senza fissare alcuna base: solo più avanti, definendo la matrice associata a $\Abb$ in una base, tradurremo questi risultati nel linguaggio dei sistemi di equazioni lineari, oggetto del capitolo successivo.


\subsection{Definizione}

Siano $X$ e $Y$ spazi vettoriali sul campo $\mathbb{R}$, di dimensione finita $n$ ed $m$ rispettivamente. Un'applicazione $\Abb: X \rightarrow Y$ si dice \textit{lineare}\index{applicazione lineare|textbf} se soddisfa le propriet\`a di additivit\`a e omogeneit\`a:
\begin{enumerate}[label=(\roman*)]
	\item $\Abb(\ub_1+\ub_2)=\Abb(\ub_1)+\Abb(\ub_2)$ \quad per tutti gli $\ub_1$, $\ub_2 \in X$,
	\item $\Abb(\alpha \ub)=\alpha \,\Abb(\ub)$ \quad per ogni $\alpha \in \mathbb{R}$ e $\ub \in X$,
\end{enumerate}
sintetizzabili nella condizione $\Abb(\alpha_1 \ub_1 + \alpha_2 \ub_2) = \alpha_1 \Abb(\ub_1) + \alpha_2 \Abb(\ub_2)$.

Fissiamo una base $(\eb_1,\dots,\eb_n)$ di $X$ e una base $(\fb_1,\dots,\fb_m)$ di $Y$. Ogni $\xb\in X$ si scrive in modo unico come $\xb=\sum_{j=1}^n x_j\eb_j$, e per linearit\`a
	\begin{equation}
		\Abb\xb=\sum_{j=1}^n x_j\,\Abb\eb_j .
	\end{equation}
	L'applicazione $\Abb$ \`e dunque completamente determinata dalle immagini $\Abb\eb_j$ dei vettori della base di $X$, espresse a loro volta nella base di $Y$:
	\begin{equation}
		\Abb\eb_j=\sum_{i=1}^m A_{ij}\,\fb_i, \qquad j=1,\dots,n.
	\end{equation}
	I coefficienti $A_{ij}$ definiscono la \textit{matrice associata} a $\Abb$ nelle basi scelte, $\bm{A}=[A_{ij}]\in\mathbb{R}^{m\times n}$: una matrice non \`e altro, dunque, che la rappresentazione di un'applicazione lineare una volta fissate le basi di dominio e codominio. In componenti, $\yb=\Abb\xb$ si scrive $y_i=\sum_j A_{ij}x_j$, ossia $\bm{y}=\bm{A}\bm{x}$, e il problema inverso $\Abb(\xb)=\bb$ diventa il sistema lineare $\bm{A}\bm{x}=\bm{b}$ che studieremo nel prossimo capitolo.



\subsection{Nucleo, immagine, rango e nullit\`a}\label{sez:nucleo_immagine}

La struttura dell'applicazione lineare $\Abb$ \`e caratterizzata da due sottospazi:
\begin{enumerate}[label=(\roman*)]
	\item il \textit{nucleo}\index{nucleo|textbf}\nomenclature[O]{$\Ker\Abb$}{nucleo di un'applicazione lineare} (o \textit{kernel})\, $\Nc(\Abb) \equiv \Ker\Abb=\{\xb\in X \mid \Abb\xb=\zerovb \}\subset X$;
	\item l'\textit{immagine}\index{immagine|textbf}\nomenclature[O]{$\Imm\Abb$}{immagine di un'applicazione lineare} (o il \textit{range})\, $\Ic(\Abb) \equiv \Imm\Abb=\{ \Abb\xb \mid \xb\in X \}\subset Y$.
\end{enumerate}
Il nucleo raccoglie tutti gli elementi che l'applicazione ``annulla''; l'immagine \`e l'insieme dei valori effettivamente raggiunti. Si definiscono il \textit{rango}\index{rango} e la \textit{nullit\`a}\index{nullità|textbf}\nomenclature[O]{$\nl\Abb$}{nullità di un'applicazione lineare} di $\Abb$:
\begin{equation}
	\rg\Abb=\dim\Ic(\Abb), \qquad \nl\Abb=\dim\Nc(\Abb).
\end{equation}

Nella rappresentazione matriciale:
	\begin{enumerate}[label=(\roman*)]
		\item $\Nc(\bm{A}) \equiv \Ker \bm{A}=\{\bm{x}\in \mathbb{R}^{n} \mid \bm{A} \bm{x}=\bm{0} \}\subset \mathbb{R}^{n}$;
		\item $\Ic(\bm{A}) \equiv \Imm\bm{A}=\Span \{\bm{a}_1,\,\bm{a}_2, \ldots, \bm{a}_n\}\subset \mathbb{R}^{m}$,
	\end{enumerate}
	dove $\bm{a}_j$ \`e la $j$-esima colonna di $\bm{A}$: per definizione stessa di $\Ic(\bm{A})$ come spazio generato dalle colonne, il rango $r:=\rg\bm{A}$ coincide con il massimo numero di colonne linearmente indipendenti di $\bm{A}$. In rango $r$ coincide anche con il massimo numero di \textit{righe} linearmente indipendenti, come dimostreremo più avanti. 


\subsection{Teorema di nullit\`a pi\`u rango}

%\begin{teorema}[Nullit\`a pi\`u rango]\label{teo:null_rango}\index{teorema!di nullità più rango|textbf}
	\emph{Per un'applicazione lineare $\Abb: X \rightarrow Y$ tra spazi di dimensione finita vale:
	\begin{equation}\label{eq:null_rango}
		\dim X = \nl\Abb + \rg\Abb.
	\end{equation}
}
%\end{teorema}

\noindent Dal teorema seguono le condizioni che legano le propriet\`a dell'applicazione alla struttura del nucleo e dell'immagine:
\begin{enumerate}[label=(\roman*)]
	\item $\Abb$ \`e \textit{iniettiva} $\;\Longleftrightarrow\;$ $\Nc(\Abb)=\{\zerovb\}$ $\;\Longleftrightarrow\;$ $\rg\Abb=\dim X =n$;
	\item $\Abb$ \`e \textit{suriettiva} $\;\Longleftrightarrow\;$ $\Ic(\Abb)=Y$ $\;\Longleftrightarrow\;$ $\rg\Abb=\dim Y = m$;
	\item se $\dim X=\dim Y$, l'iniettivit\`a equivale alla suriettivit\`a (e alla biettivit\`a).
\end{enumerate}

\noindent Nella rappresentazione matriciale, per $\bm{A}\in\mathbb{R}^{m\times n}$ di rango $r$, il teorema fornisce le dimensioni esplicite dei sottospazi:
\begin{equation}
	\dim\Nc(\bm{A}) = n - r, \qquad \dim\Ic(\bm{A}) = r.
\end{equation}
Il nucleo \`e dunque un sottospazio di $\mathbb{R}^n$ di dimensione $n-r$ (banale se e solo se $r=n$), mentre l'immagine \`e un sottospazio di $\mathbb{R}^m$ di dimensione $r$ (coincidente con $\mathbb{R}^m$ se e solo se $r=m$).

\begin{oss}
	La linearit\`a trasforma le questioni di esistenza e unicit\`a del problema inverso (Sez.~\ref{par:problema_inverso}) in questioni sulla dimensione di due sottospazi: il problema inverso $\bm{A}\bm{x}=\bm{b}$ ammette soluzione se e solo se $\bm{b} \in \Ic(\bm{A})$ (esistenza), e la soluzione \`e unica se e solo se $\Nc(\bm{A})=\{\bm{0}\}$ (unicit\`a). Si noti la differenza rispetto al caso generale non lineare (\SECref{par:problema_inverso}): l\`i, l'iniettivit\`a di $f$ garantiva l'unicit\`a per \textit{ogni} $\bar y$, ma non era necessaria per l'unicit\`a relativa a un singolo $\bar y$ fissato (l'esempio $f(x)=x^2$ in $\bar y=0$). Qui invece, per la linearit\`a, se $\xb_0$ e $\xb_1$ risolvono entrambi $\Abb\xb=\bb$, la loro differenza $\xb_1-\xb_0$ appartiene al nucleo: tutte le controimmagini non vuote di $\Abb$ hanno dunque la stessa struttura, quella di $\Nc(\Abb)$, e l'equivalenza fra iniettivit\`a e unicit\`a torna a valere anche per un singolo $\bar y$ fissato. Quando il nucleo \`e non banale, le soluzioni --- se esistono --- differiscono tra loro per elementi del nucleo, come espresso dalla struttura $\bm{x} = \bm{x}_0 + \bm{X}_\alpha\bm{\alpha}$, che ritroveremo nel prossimo capitolo nello studio dei sistemi lineari.
\end{oss}


%-------------------------------------------------------------------------------------------------------------------
\section{Operatore aggiunto e sottospazi fondamentali}
%-------------------------------------------------------------------------------------------------------------------

Abbiamo visto che il problema $\Abb\xb=\bb$ ammette soluzione se e solo se $\bb\in\Ic(\Abb)$: una condizione corretta, ma che di fatto ripete la domanda invece di rispondervi, dato che verificare l'appartenenza a $\Ic(\Abb)$ equivale a provare a risolvere il problema stesso. Conviene allora affiancare ad $\Abb: X \rightarrow Y$ un secondo operatore che agisca in direzione opposta, da $Y$ verso $X$: vedremo che questo operatore, detto \textit{aggiunto} di $\Abb$, permette di riscrivere la condizione di esistenza in una forma verificabile in modo indipendente, e che i quattro sottospazi coinvolti --- nucleo e immagine di $\Abb$, nucleo e immagine del suo aggiunto --- sono legati da relazioni di ortogonalità che completano il quadro geometrico del problema. Lo stesso operatore aggiunto fornisce, più avanti in questo capitolo, il fondamento algebrico del \textit{problema duale}: nella meccanica delle strutture, come vedremo, è ciò che lega la descrizione cinematica di una struttura alla sua descrizione statica.

\subsection{Operatore aggiunto}\label{sec:aggiunto_sottospazi}

Se gli spazi $X$ e $Y$ sono dotati di prodotto scalare, per ogni $\yb\in Y$ fissato l'applicazione $\xb\mapsto\Abb\xb\cdot\yb$ \`e  lineare su $X$ (lineare perch\'e lo sono $\Abb$ e il prodotto scalare). Per il risultato di rappresentazione visto nella  \SECref{prodvet}  (Teorema di Riesz, v. nota \ref{riesz}) --- che vale per  un'applicazione lineare a valori reali  su un qualunque spazio vettoriale di dimensione finita con prodotto scalare, non solo su $\Uc$ --- esiste un unico $\zb\in X$ tale che $\xb\cdot\zb=\Abb\xb\cdot\yb$ per ogni $\xb\in X$; presa una base ortonormale $\{\eb_1,\dots,\eb_n\}$ di $X$, tale $\zb$ \`e dato esplicitamente da $\zb=\sum_{j=1}^n(\Abb\eb_j\cdot\yb)\,\eb_j$. Poniamo $\Abb^\ast\yb:=\zb$.

Questo fissa $\Abb^\ast\yb$ per ogni singolo $\yb$. Verifichiamo che $\Abb^\ast$ sia un'applicazione lineare. Per $\yb_1,\yb_2\in Y$ e $\alpha_1,\alpha_2\in\mathbb{R}$, e per ogni $\xb\in X$,
\[
\begin{aligned}
	\xb\cdot\Abb^\ast(\alpha_1\yb_1+\alpha_2\yb_2)&=\Abb\xb\cdot(\alpha_1\yb_1+\alpha_2\yb_2)=\alpha_1(\xb\cdot\Abb^\ast\yb_1)+\alpha_2(\xb\cdot\Abb^\ast\yb_2)\\
	&=\xb\cdot(\alpha_1\Abb^\ast\yb_1+\alpha_2\Abb^\ast\yb_2)
\end{aligned}
\]
per ogni $\xb\in X$; per l'unicit\`a nel risultato di rappresentazione (applicata a $\xb=\Abb^\ast(\alpha_1\yb_1+\alpha_2\yb_2)-(\alpha_1\Abb^\ast\yb_1+\alpha_2\Abb^\ast\yb_2)$), questo forza $\Abb^\ast(\alpha_1\yb_1+\alpha_2\yb_2)=\alpha_1\Abb^\ast\yb_1+\alpha_2\Abb^\ast\yb_2$. Dunque $\Abb^\ast:Y\to X$ \`e lineare: lo chiamiamo l'\textit{aggiunto}\index{operatore aggiunto|textbf}\nomenclature[O]{$\Abb^{\ast}$}{operatore aggiunto} di $\Abb$, l'unica applicazione lineare tale che
\begin{equation}\label{eq:aggiunto}
	\Abb\xb\cdot\yb=\xb\cdot\Abb^\ast\yb
	\qquad \forall\,\xb \in X,\;\forall\,\yb \in Y.
\end{equation}

Nella rappresentazione matriciale, se le basi scelte per $X$ e $Y$ sono ortonormali --- come sar\`a sempre il caso nel seguito --- la matrice associata ad $\Abb^\ast$ \`e la trasposta $\bm{A}^\top \in \mathbb{R}^{n\times m}$. Per questo motivo, scriveremo anche $\Abb^\to$ piuttosto che $\Abb^\ast$


All'operatore aggiunto sono associati i sottospazi:
\begin{enumerate}[label=(\roman*)]
	\item $\Nc(\Abb^\top) \equiv \Ker \Abb^\top=\{\yb\in Y \mid \Abb^\top \yb=\zerovb \}\subset Y$;
	\item $\Ic(\Abb^\top) \equiv \Imm\Abb^\top=\Span \{\bm{a}^1,\,\bm{a}^2, \ldots, \bm{a}^m\}\subset X$,
\end{enumerate}
dove $\bm{a}^i$ \`e la $i$-esima riga di $\bm{A}$. Per il teorema di nullit\`a pi\`u rango applicato ad $\Abb^\top:Y\to X$ vale
\begin{equation}
	\dim Y=\nl\Abb^\top+\rg\Abb^\top.
\end{equation}
Vedremo fra poco che $\rg\Abb^\top=\rg\Abb$: non \`e un fatto immediato, ma conseguenza delle relazioni di ortogonalit\`a fra i sottospazi fondamentali che introduciamo ora --- ed \`e proprio ci\`o che chiude la questione, lasciata in sospeso, dell'uguaglianza fra numero massimo di colonne e di righe indipendenti di $\bm{A}$.



\subsection{Relazioni di ortogonalit\`a}\index{sottospazi fondamentali|textbf}

I quattro sottospazi $\Nc(\Abb)$, $\Ic(\Abb)$, $\Nc(\Abb^\top)$, $\Ic(\Abb^\top)$ sono legati dalle seguenti relazioni fondamentali:\footnote{Dimostriamo che $\Nc(\Abb) = \bigl(\Ic(\Abb^\top)\bigr)^\perp$. 
	Se $\xb \in \Nc(\Abb)$, allora $\Abb\xb = \zerovb$ e per ogni 
	$\yb \in Y$ si ha $\xb \cdot \Abb^\top\yb = \Abb\xb \cdot \yb = 
	\zerovb \cdot \yb = 0$; dunque $\xb$ \`e ortogonale a ogni elemento 
	di $\Ic(\Abb^\top)$, cio\`e $\xb \in \bigl(\Ic(\Abb^\top)\bigr)^\perp$. 
	Viceversa, se $\xb \in \bigl(\Ic(\Abb^\top)\bigr)^\perp$, allora 
	$\xb \cdot \Abb^\top\yb = 0$ per ogni $\yb$, cio\`e 
	$\Abb\xb \cdot \yb = 0$ per ogni $\yb \in Y$; in particolare per $\yb=\Abb\xb$, da cui $|\Abb\xb|^2=0$, ossia
	$\Abb\xb = \zerovb$, cio\`e $\xb \in \Nc(\Abb)$. 
	L'altra relazione, $\Nc(\Abb^\top) = \bigl(\Ic(\Abb)\bigr)^\perp$, 
	si dimostra in modo del tutto analogo scambiando i ruoli di $\Abb$ 
	e $\Abb^\top$: essendo il prodotto scalare simmetrico, la relazione $\Abb\xb\cdot\yb=\xb\cdot\Abb^\top\yb$ si legge infatti anche come $\Abb^\top\yb\cdot\xb=\yb\cdot\Abb\xb$ per ogni $\xb,\yb$, cio\`e come la relazione che definisce l'aggiunto di $\Abb^\top$ --- che \`e dunque $\Abb$ stesso.}
\begin{equation}\label{eq:ortogonalita}
	\Nc(\Abb^\top)=\bigl(\Ic(\Abb)\bigr)^\perp, \qquad
	\Nc(\Abb)=\bigl(\Ic(\Abb^\top)\bigr)^\perp,
\end{equation}
dove $W^\perp = \{\xb \in X \mid \xb \cdot \wb = 0 \;\; \forall\, \wb \in W\}$ denota il complemento ortogonale\index{complemento ortogonale|textbf}\nomenclature[O]{$W^{\perp}$}{complemento ortogonale di un sottospazio} di un sottospazio $W$. Ne segue che ciascuno dei due spazi $X$ e $Y$ si decompone come somma diretta di sottospazi ortogonali:
\begin{equation}\label{eq:somma_diretta}
	X=\Nc(\Abb) \oplus \Ic(\Abb^\top), \qquad Y=\Nc(\Abb^\top) \oplus \Ic(\Abb).
\end{equation}
%---------------------------------------------------------------------------------------------------------------------
\begin{figure}[htbp]
	\centering
	\begin{tikzpicture}[>=Stealth, line width=0.8pt, scale=0.8]
		\draw[thick] (-3.5, 0) ellipse [x radius=2.2, y radius=2.8];
		\node at (-3.5, 3.2) {$\mathbb{R}^n$};
		\draw[thick] (-5.69, 0.3)
		.. controls (-4.5, 0.6) and (-2.5, -0.6) ..
		(-1.31, -0.3);
		\node at (-4.3,  1.6) {$\Nc(\Abb)$};
		\fill (-4.0,  0.9) circle (1.5pt) node[right=2pt] {$\xb_n$};
		\node at (-2.7, -1.6) {$\Ic(\Abb^\top)$};
		\fill (-3.0, -0.9) circle (1.5pt) node[right=2pt] {$\xb_c$};
		\draw[thick] (3.5, 0) ellipse [x radius=2.2, y radius=2.8];
		\node at (3.5, 3.2) {$\mathbb{R}^m$};
		\draw[thick] (1.31, 0.3)
		.. controls (2.5, 0.6) and (4.5, -0.6) ..
		(5.69, -0.3);
		\node at (2.7,  1.6) {$\Ic(\Abb)$};
		\fill (3.0,  0.9) circle (1.5pt) node[right=2pt] {$\yb_c$};
		\node at (4.3, -1.6) {$\Nc(\Abb^\top)$};
		\fill (4.0, -0.9) circle (1.5pt) node[right=2pt] {$\yb_n$};
		\draw[->, thick] (-1.0, 1.5) to[bend left=20]
		node[midway, above=4pt] {$\Abb_{m\times n}$}
		(1.0, 1.5);
		\draw[->, thick] (1.0, -1.5) to[bend left=20]
		node[midway, below=4pt] {$\Abb^\top_{n\times m}$}
		(-1.0, -1.5);
	\end{tikzpicture}
	\caption{I quattro sottospazi fondamentali di $\Abb_{m\times n}$.
		Ogni $\xb \in \mathbb{R}^n$ si decompone in modo unico come
		$\xb = \xb_n + \xb_c$ con $\xb_n \in \Nc(\Abb)$ (componente
		nel \emph{nucleo}) e $\xb_c \in \Ic(\Abb^\top)$ (componente
		\emph{complementare} al nucleo); analogamente ogni
		$\yb \in \mathbb{R}^m$ si decompone come $\yb = \yb_c + \yb_n$
		con $\yb_c \in \Ic(\Abb)$ e $\yb_n \in \Nc(\Abb^\top)$.
		L'operatore $\Abb$ mappa $\xb_c$ in $\yb_c$ e annulla $\xb_n$;
		l'operatore $\Abb^\top$ mappa $\yb_c$ in $\xb_c$ e annulla $\yb_n$.}
	\label{fig:quattro_sottospazi}
\end{figure}
%---------------------------------------------------------------------------------------------------------------------

Le dimensioni dei quattro sottospazi sono vincolate dal teorema di nullit\`a pi\`u rango, applicato sia ad $\Abb$ sia al suo aggiunto $\Abb^\top$, combinato con le relazioni di ortogonalit\`a appena dimostrate. Poniamo $r:=\rg\Abb=\dim\Ic(\Abb)$: il teorema, applicato ad $\Abb$, d\`a $\dim\Nc(\Abb)=n-r$. Per la relazione $\Nc(\Abb^\top)=\bigl(\Ic(\Abb)\bigr)^\perp$, e poich\'e il complemento ortogonale di un sottospazio di dimensione $r$ in uno spazio di dimensione $m$ ha dimensione $m-r$, si ha $\dim\Nc(\Abb^\top)=m-r$. Applicando allora il teorema di nullit\`a pi\`u rango ad $\Abb^\top:Y\to X$:
\begin{equation}
	\dim Y=\dim\Nc(\Abb^\top)+\dim\Ic(\Abb^\top) \;\Longrightarrow\; \dim\Ic(\Abb^\top)=m-(m-r)=r.
\end{equation}
Abbiamo dunque dimostrato $\rg\Abb^\top=\dim\Ic(\Abb^\top)=r=\rg\Abb$: il rango dell'aggiunto coincide con quello di $\Abb$. Questo chiude anche la questione lasciata in sospeso a proposito della matrice $\bm A$ (\SECref{sez:nucleo_immagine}): poich\'e $\Ic(\bm A^\top)$ \`e generato dalle righe di $\bm A$, l'uguaglianza $\rg\bm A^\top=\rg\bm A$ dice esattamente che il massimo numero di righe linearmente indipendenti coincide con quello delle colonne.

Riassumendo, le dimensioni dei quattro sottospazi fondamentali sono:
\begin{equation}\label{eq:dim_4spazi}
	\underbrace{\dim\Nc(\Abb)}_{n-r} + \underbrace{\dim\Ic(\Abb^\top)}_{r} = n, \qquad
	\underbrace{\dim\Nc(\Abb^\top)}_{m-r} + \underbrace{\dim\Ic(\Abb)}_{r} = m.
\end{equation}




\paragraph{Condizione di esistenza.}\label{sec:rilettura_classificazione}
Le relazioni di ortogonalità permettono di riformulare in termini puramente geometrici le condizioni di esistenza e unicità del problema $\Abb\xb=\bb$ in \SECref{par:problema_inverso}. Poiché $\Ic(\Abb)=\bigl(\Nc(\Abb^\top)\bigr)^\perp$, l'esistenza della soluzione equivale a
\begin{equation}\label{eq:compatibilita}
	\yb\cdot\bb=0 \qquad \forall\,\yb\in\Nc(\Abb^\top):
\end{equation}
$\bb$ deve essere ortogonale a ogni elemento del nucleo dell'aggiunto.

\paragraph{Struttura della soluzione.}
L'unicità, come già visto, equivale a $\Nc(\Abb)=\{\zerovb\}$; quando il nucleo non è banale, le soluzioni --- se esistono --- differiscono fra loro per un elemento di $\Nc(\Abb)$, e formano una famiglia affine di dimensione $\dim\Nc(\Abb)=n-r$.



%\begin{oss}
%	Nella meccanica delle strutture, il problema $\Abb\xb=\bb$ descriverà un problema cinematico o un problema statico. Il nucleo $\Nc(\Abb)$ rappresenterà, a seconda del caso, i meccanismi cinematici o le autotensioni; il nucleo dell'aggiunto $\Nc(\Abb^\top)$ rappresenterà i carichi non equilibrabili o le deformazioni incompatibili. La condizione di ortogonalità \eqref{eq:compatibilita} tradurrà, in linguaggio algebrico, principi fisici fondamentali quali l'equilibrio e la congruenza. Questa corrispondenza fra struttura algebrica e significato fisico costituirà il filo conduttore dell'intero testo.
%\end{oss}
%
%I quattro sottospazi fondamentali non esauriscono però il loro ruolo nella classificazione del problema $\Abb\xb=\bb$. Essi costituiscono lo strumento naturale per comprendere la struttura di un secondo problema, detto \textit{duale}, che affianca simmetricamente il primale e con esso condivide una identità fondamentale --- l'identità bilineare --- che nelle applicazioni fisiche assumerà il significato del principio dei lavori virtuali. La classificazione completa, estesa al problema duale, è presentata nella Sez.~\ref{sec:risolvibilita_duale}.
%	

%-------------------------------------------------------------------------------------------------------------------
	\section{Il problema duale}
	%-------------------------------------------------------------------------------------------------------------------
	
	Le relazioni di ortogonalità appena viste mostrano che l'aggiunto $\Abb^\top$ non è solo uno strumento ausiliario per verificare la condizione di esistenza del problema $\Abb\xb=\bb$: è un'applicazione lineare $\Abb^\top:Y\to X$ a pieno titolo, tanto quanto lo è $\Abb$. Ha dunque senso porre, per $\Abb^\top$, lo stesso problema che abbiamo posto per $\Abb$: dato $\cb\in X$, trovare $\yb\in Y$ tale che $\Abb^\top\yb=\cb$. Chiamiamo questo secondo problema \textit{duale} del primo (che chiamiamo, per contrasto, \textit{primale}). La denominazione non è arbitraria: avendo già mostrato che l'aggiunto di $\Abb^\top$ è $\Abb$ stesso, la relazione fra i due problemi è perfettamente simmetrica --- non c'è un problema principale e uno accessorio, ma una coppia di problemi sullo stesso piano, ciascuno aggiunto dell'altro. Nella meccanica delle strutture, come vedremo, è proprio questa coppia a descrivere il legame fra il problema cinematico e quello statico di una struttura.

\subsection{Formulazione}

Il problema $\Abb\xb=\bb$, già introdotto in \SECref{par:problema_inverso}, si chiama in questo contesto problema \textit{primale}\index{problema primale|textbf}; il problema $\Abb^\top\yb=\cb$ appena introdotto si chiama problema \textit{duale}\index{problema duale|textbf}:
\begin{equation}\label{eq:primale}
	\Abb\xb=\bb \qquad \text{(primale)}, \qquad\qquad
	\Abb^\top\yb=\cb \qquad \text{(duale)},
\end{equation}
con $\xb\in X$, $\bb\in Y$ per il primale, e $\yb\in Y$, $\cb\in X$ per il duale. Si noti dove vivono le quattro grandezze: la variabile duale $\yb$ appartiene allo stesso spazio $Y$ del termine noto del primale $\bb$; il termine noto duale $\cb$ appartiene allo stesso spazio $X$ della variabile primale $\xb$. 

Nella rappresentazione matriciale, fissate basi ortonormali di $X$ e $Y$, i due problemi diventano i sistemi lineari
\begin{equation}\label{eq:duale}
	\bm{A}\,\bm{x} = \bm{b} \qquad \text{(primale)}, \qquad\qquad
	\bm{A}^\top\,\bm{y} = \bm{c} \qquad \text{(duale)},
\end{equation}
con $\bm{A} \in \mathbb{R}^{m \times n}$, $\bm{x} \in \mathbb{R}^n$, $\bm{b} \in \mathbb{R}^m$, $\bm{y} \in \mathbb{R}^m$ vettore delle variabili duali, $\bm{c} \in \mathbb{R}^n$ vettore dei termini noti duali.



\subsection{Interpretazione mediante combinazioni lineari}

La natura dei due problemi si chiarisce riscrivendoli in forma di
combinazioni lineari:
\begin{equation}\label{eq:comb_lin_duale}
	\sum_{j=1}^n x_j\,\bm{a}_j = \bm{b}, \qquad
	\sum_{i=1}^m y_i\,\bm{a}^i = \bm{c},
\end{equation}
dove $\bm{a}_j \in \mathbb{R}^m$ \`e la $j$-esima colonna di
$\bm{A}$ e $\bm{a}^i \in \mathbb{R}^n$ \`e la $i$-esima riga di
$\bm{A}$ scritta come vettore colonna --- ovvero la $i$-esima
colonna di $\bm{A}^\top$. Risolvere il problema primale significa
trovare i coefficienti $x_j$ che esprimono $\bm{b}$ come
combinazione lineare delle \textit{colonne} di $\bm{A}$; risolvere
il problema duale significa trovare i coefficienti $y_i$ che
esprimono $\bm{c}$ come combinazione lineare delle \textit{colonne}
di $\bm{A}^\top$ (cio\`e delle righe di $\bm{A}$, intese come
vettori colonna).


\subsection{Rappresentazione schematica della dualit\`a}\index{dualità!algebra lineare|textbf}

La struttura simmetrica dei due problemi \`e sintetizzata nello schema seguente, che si legge per righe (problema primale) e per colonne (problema duale):

\begin{figure}[H]%[ht]
	\centering
	\begin{equation*}
		\begin{array}{c|cccccc|c}
			& x_1 & x_2 & \dots & x_j & \dots & x_n & \\
			\hline
			y_1 & a_{11} & a_{12} & \dots & a_{1j} & \dots & a_{1n} & b_1 \\
			y_2 & a_{21} & a_{22} & \dots & a_{2j} & \dots & a_{2n} & b_2 \\
			\vdots & \vdots & \vdots & \ddots & \vdots & \ddots & \vdots & \vdots \\
			y_i & a_{i1} & a_{i2} & \dots & a_{ij} & \dots & a_{in} & b_i \\
			\vdots & \vdots & \vdots & \ddots & \vdots & \ddots & \vdots & \vdots \\
			y_m & a_{m1} & a_{m2} & \dots & a_{mj} & \dots & a_{mn} & b_m \\
			\hline
			& c_1 & c_2 & \dots & c_j & \dots & c_n &
		\end{array}
	\end{equation*}
	\caption{Rappresentazione schematica della dualit\`a nei sistemi lineari.}
	\label{fig:duale}
\end{figure}

\noindent La dualit\`a si manifesta nelle seguenti corrispondenze:
\begin{enumerate}[label=(\roman*)]
	\item ogni variabile $x_j$ del problema primale corrisponde a un'equazione (la $j$-esima) del problema duale;
	\item ogni variabile $y_i$ del problema duale corrisponde a un'equazione (la $i$-esima) del problema primale;
	\item i due problemi sono governati da matrici trasposte l'una dell'altra.
\end{enumerate}


\begin{oss}[\textbf{Reciprocità fra primale e duale}]
	Le soluzioni di primale e duale sono legate da una relazione immediata: se $\Abb\xb=\bb$ e $\Abb^\top\yb=\cb$, allora
	\begin{equation}\label{eq:id_bil}
		\yb\cdot\bb=\Abb\xb\cdot\yb=\xb\cdot\Abb^\top\yb=\xb\cdot\cb,
	\end{equation}
	dove il passaggio centrale usa solo la definizione di aggiunto \eqref{eq:aggiunto}. Non è un'equazione da risolvere né una condizione da verificare, ma una conseguenza automatica: chiunque disponga di una soluzione del primale e di una del duale la ottiene gratuitamente. Nella rappresentazione matriciale, $\bm{y}^\top\bm{b}=\bm{x}^\top\bm{c}$. 
\end{oss}






%-------------------------------------------------------------------------------------------------------------------
\section{Risolvibilità e sottospazi fondamentali}	\label{sec:risolvibilita_duale}
%-------------------------------------------------------------------------------------------------------------------

La condizione di esistenza della soluzione, espressa per il solo
problema primale nel \PARref{sec:rilettura_classificazione}, si estende
simmetricamente al problema duale: il primale $\Abb\xb=\bb$
ammette soluzione se e solo se $\bb\perp\Nc(\Abb^\top)$; il
duale $\Abb^\top\yb=\cb$ ammette soluzione se e solo se
$\cb\perp\Nc(\Abb)$ (per la relazione $\Ic(\Abb^\top)=\bigl(\Nc(\Abb)\bigr)^\perp$, conseguenza di \eqref{eq:ortogonalita}). La simmetria è completa: l'ostacolo
all'esistenza risiede sempre nel nucleo dell'operatore opposto.

Nella rappresentazione matriciale, $\bm{A}\bm{x}=\bm{b}$ ammette soluzione se e solo se $\bm{b}\perp\Nc(\bm{A}^\top)$, e $\bm{A}^\top\bm{y}=\bm{c}$ se e solo se $\bm{c}\perp\Nc(\bm{A})$.

\begin{oss}
	La \FIGref{fig:problemi_duali} illustra la struttura comune ai due
	problemi. Ogni elemento di $\mathbb{R}^n$ si decompone in una
	componente \emph{complementare al nucleo} $\bm{x}_c \in \Ic(\bm{A}^\top)$
	e una componente nel nucleo $\bm{x}_n \in \Nc(\bm{A})$ --- dove il
	pedice $c$ sta per \emph{complementare} e il pedice $n$ per
	\emph{nucleo}. Analogamente ogni elemento di $\mathbb{R}^m$ si
	decompone in $\bm{y}_c \in \Ic(\bm{A})$ e $\bm{y}_n \in
	\Nc(\bm{A}^\top)$. Le condizioni di compatibilità $\bm{b}_n = \bm{0}$
	e $\bm{c}_n = \bm{0}$ esprimono il fatto che il dato deve appartenere
	all'immagine del rispettivo operatore: la sua componente nel nucleo
	dell'operatore duale deve essere nulla. La soluzione generale è
	determinata a meno di elementi arbitrari del nucleo:
	$\bm{x} = \bm{x}_c + \bm{x}_n$ con $\bm{x}_n \in \Nc(\bm{A})$
	libero, e $\bm{y} = \bm{y}_c + \bm{y}_n$ con
	$\bm{y}_n \in \Nc(\bm{A}^\top)$ libero.
\end{oss}
%---------------------------------------------------------------------------------------------------------------------
\begin{figure}[htbp]
	\centering
	\begin{tikzpicture}[>=Stealth, line width=0.8pt, scale=0.8]
		\draw[thick] (-3.5, 0) ellipse [x radius=2.2, y radius=2.8];
		\node at (-3.5, 3.2) {$\mathbb{R}^n$};
		\draw[thick] (-5.7, 0)
		.. controls (-4.5, 0.5) and (-2.5, -0.5) ..
		(-1.3, 0);
		\node at (-3.5,  2.0) {$\Nc(\bm{A})$};
		\fill (-3.5,  1.2) circle (1.5pt) node[right=2pt] {$\bm{x}_n$};
		\node[gray!80] at (-3.5,  0.5) {$\bm{c}_n = \bm{0}$};
		\fill (-3.5, -0.5) circle (1.5pt) node[right=2pt] {$\bm{x}_c$};
		\fill (-3.5, -1.2) circle (1.5pt) node[right=2pt] {$\bm{c}_c$};
		\node at (-3.5, -2.0) {$\Ic(\bm{A}^\top)$};
		\draw[thick] (3.5, 0) ellipse [x radius=2.2, y radius=2.8];
		\node at (3.5, 3.2) {$\mathbb{R}^m$};
		\draw[thick] (1.3, 0)
		.. controls (2.5, 0.5) and (4.5, -0.5) ..
		(5.7, 0);
		\node at (3.5,  2.0) {$\Ic(\bm{A})$};
		\fill (3.5,  1.2) circle (1.5pt) node[right=2pt] {$\bm{b}_c$};
		\fill (3.5,  0.5) circle (1.5pt) node[right=2pt] {$\bm{y}_c$};
		\node[gray!80] at (3.5, -0.5) {$\bm{b}_n = \bm{0}$};
		\fill (3.5, -1.2) circle (1.5pt) node[right=2pt] {$\bm{y}_n$};
		\node at (3.5, -2.0) {$\Nc(\bm{A}^\top)$};
		\draw[->, thick] (-1.0, 1.5) to[bend left=20]
		node[midway, above=4pt] {$\bm{A}\bm{x}_c = \bm{b}_c$}
		(1.0, 1.5);
		\draw[->, thick] (1.0, -1.5) to[bend left=20]
		node[midway, below=4pt] {$\bm{A}^\top\bm{y}_c = \bm{c}_c$}
		(-1.0, -1.5);
	\end{tikzpicture}
	\caption{Struttura dei problemi primale $\bm{A}\bm{x} = \bm{b}$ e duale
		$\bm{A}^\top\bm{y} = \bm{c}$. Il pedice $c$ denota la componente
		\emph{complementare al nucleo}: $\bm{x}_c \in \Ic(\bm{A}^\top)$,
		$\bm{b}_c \in \Ic(\bm{A})$, $\bm{y}_c \in \Ic(\bm{A})$,
		$\bm{c}_c \in \Ic(\bm{A}^\top)$. Il pedice $n$ denota la componente
		nel nucleo: $\bm{x}_n \in \Nc(\bm{A})$, $\bm{y}_n \in \Nc(\bm{A}^\top)$.
		Il primale ammette soluzione se e solo se $\bm{b}_n = \bm{0}$;
		il duale se e solo se $\bm{c}_n = \bm{0}$.}
	\label{fig:problemi_duali}
\end{figure}
%---------------------------------------------------------------------------------------------------------------------

\subsection{Classificazione dei problemi primale e duale}
\label{sec:classificazione_primale_duale}

La tabella seguente raccoglie la classificazione completa in termini
dei quattro sottospazi fondamentali, estesa a entrambi i problemi.


\begin{table}[H]%[ht]
	\centering
	\small
	\renewcommand{\arraystretch}{1.4}
	\begin{tabular}{@{}lcccc@{}}
		\toprule
		\textbf{Tipo} & $\dim\Nc(\bm{A})$ & $\dim\Nc(\bm{A}^\top)$ & \textbf{Primale} & \textbf{Duale} \\
		\midrule
		Isodeterminato   & $0$   & $0$   & unica $\forall\,\bm{b}$ & unica $\forall\,\bm{c}$ \\
		Sottodeterminato & $n{-}m$ & $0$   & $\infty^{n-m}\;\forall\,\bm{b}$ & unica se $\bm{c}\perp\Nc(\bm{A})$ \\
		Sovradeterminato & $0$   & $m{-}n$ & unica se $\bm{b}\perp\Nc(\bm{A}^\top)$ & $\infty^{m-n}\;\forall\,\bm{c}$ \\
		Degenere         & $n{-}r$ & $m{-}r$ & $\infty^{n-r}$ se compat. & $\infty^{m-r}$ se compat. \\
		\bottomrule
	\end{tabular}
	\caption{Classificazione dei problemi primale e duale in termini dei sottospazi fondamentali.}
	\label{tab:primale_duale}
\end{table}

\begin{oss}
	La tabella mostra una simmetria notevole: un sistema primale sottodeterminato ha un duale sovradeterminato, e viceversa. I sistemi isodeterminati sono autoduali. I sistemi degeneri restano degeneri anche nel duale, ma con i ruoli di $n-r$ e $m-r$ scambiati. In ogni caso, la relazione di reciprocità $\bm{y}^\top\bm{b} = \bm{x}^\top\bm{c}$ resta valida. Questa struttura duale riapparirà sistematicamente nel corso, dove il problema cinematico e il problema statico di una struttura saranno l'uno il duale dell'altro.
\end{oss}


\begin{oss}
	La classificazione della Tab.~\ref{tab:primale_duale} determina non
	solo l'esistenza della soluzione ma anche la sua forma esplicita: in
	tutti i casi la soluzione generale si scrive come somma di una
	soluzione particolare e della combinazione lineare più generale di
	elementi del nucleo, con la stessa forma per entrambi i problemi e i
	ruoli di $\Abb$ e $\Abb^\top$ scambiati. Le formule esplicite,
	che richiedono di saper calcolare l'inversa di una matrice e una
	base del nucleo, sono sviluppate nel capitolo successivo, dove
	la classificazione qui ottenuta sarà tradotta caso per caso in
	soluzioni esplicite.
\end{oss}


%
%
%\todo[inline]{Eliminerei gli esercizi}
%
%%-------------------------------------------------------------------------------------------------------------------
%\section{Esempi numerici: primale e duale a confronto}
%%-------------------------------------------------------------------------------------------------------------------
%
%Si riprendono gli esempi dell'\APPref{app:algebra_eqlineari}, affiancando a ciascun sistema primale il corrispondente problema duale. Per ogni caso si identificano i quattro sottospazi fondamentali e si verifica l'identit\`a bilineare.
%
%\subsection{Sistema isodeterminato ($m=n=r=2$)}
%
%\textit{Primale:}\quad
%$\bm{A}\bm{x}=\bm{b}$\; con\;
%$\bm{A}=\bigl[\begin{smallmatrix}2 & 1\\ 1 & 1\end{smallmatrix}\bigr]$,\;
%$\bm{b}=(3,\;2)^\top$.\quad
%Soluzione: $\bm{x}=(1,\;1)^\top$.
%
%\medskip\noindent
%\textit{Duale:}\quad
%$\bm{A}^\top\bm{y}=\bm{c}$\; con\;
%$\bm{A}^\top=\bigl[\begin{smallmatrix}2 & 1\\ 1 & 1\end{smallmatrix}\bigr]$,\;
%$\bm{c}=(1,\;3)^\top$.\quad
%Essendo $\bm{A}$ simmetrica, il duale ha la stessa matrice. La soluzione unica \`e:
%\begin{equation*}
%	\bm{y} = (\bm{A}^\top)^{-1}\bm{c}
%	= \begin{bmatrix} 1 & -1 \\ -1 & 2 \end{bmatrix}
%	\begin{Bmatrix} 1 \\ 3 \end{Bmatrix}
%	= \begin{Bmatrix} -2 \\ 5 \end{Bmatrix}.
%\end{equation*}
%
%\noindent\textit{Identit\`a bilineare:}
%\begin{equation*}
%	\bm{y}^\top\bm{b} = (-2)(3)+(5)(2) = 4, \qquad
%	\bm{x}^\top\bm{c} = (1)(1)+(1)(3) = 4. \quad\checkmark
%\end{equation*}
%
%\noindent\textit{Sottospazi:}\quad
%$\Nc(\bm{A})=\{\bm{0}\}$,\; $\Ic(\bm{A})=\mathbb{R}^2$,\;
%$\Nc(\bm{A}^\top)=\{\bm{0}\}$,\; $\Ic(\bm{A}^\top)=\mathbb{R}^2$.\;
%Entrambi i problemi ammettono soluzione unica per qualunque termine noto.
%
%
%\subsection{Sistema sottodeterminato ($m=r=1$, $n=2$)}
%
%\textit{Primale:}\quad
%$\bm{A}\bm{x}=\bm{b}$\; con\;
%$\bm{A}=\bigl[1\;\;2\bigr]$,\;
%$b=4$.\quad
%Soluzione: $\bm{x}=(4,\;0)^\top+(-2,\;1)^\top\alpha$, con $\alpha$ libero.
%
%\medskip\noindent
%\textit{Duale:}\quad
%$\bm{A}^\top y=\bm{c}$\; con\;
%$\bm{A}^\top=\bigl(\begin{smallmatrix}1\\ 2\end{smallmatrix}\bigr)$,\;
%$\bm{c}=(c_1,\;c_2)^\top$.
%
%\noindent Il duale \`e un sistema sovradeterminato ($2$ equazioni, $1$ incognita):
%\begin{equation*}
%	\begin{Bmatrix} 1 \\ 2 \end{Bmatrix} y
%	= \begin{Bmatrix} c_1 \\ c_2 \end{Bmatrix}.
%\end{equation*}
%La condizione di compatibilit\`a \`e $c_2 = 2\,c_1$; quando \`e soddisfatta,
%la soluzione unica \`e $y=c_1$.
%
%\medskip\noindent
%\textit{Verifica dell'identit\`a bilineare} (con $\alpha=1$, cio\`e
%$\bm{x}=(2,\;1)^\top$, e $\bm{c}=(1,\;2)^\top$, compatibile, $y=1$):
%\begin{equation*}
%	\bm{y}^\top\bm{b} = (1)(4) = 4, \qquad
%	\bm{x}^\top\bm{c} = (2)(1)+(1)(2) = 4. \quad\checkmark
%\end{equation*}
%
%\noindent\textit{Sottospazi:}
%\begin{equation*}
%	\begin{aligned}
%		&\Nc(\bm{A}) = \Span\{(-2,\;1)^\top\} \subset \mathbb{R}^2, \quad
%		&&\Ic(\bm{A}) = \mathbb{R}^1, \\
%		&\Nc(\bm{A}^\top) = \{0\} \subset \mathbb{R}^1, \quad
%		&&\Ic(\bm{A}^\top) = \Span\{(1,\;2)^\top\} \subset \mathbb{R}^2.
%	\end{aligned}
%\end{equation*}
%Si verifica l'ortogonalit\`a: $(-2,\;1)\cdot(1,\;2)=-2+2=0$, cio\`e
%$\Nc(\bm{A}) \perp \Ic(\bm{A}^\top)$, e la somma diretta
%$\mathbb{R}^2 = \Nc(\bm{A}) \oplus \Ic(\bm{A}^\top)$.
%
%\begin{oss}
%	La simmetria della dualit\`a \`e evidente: il primale sottodeterminato
%	($\infty^1$ soluzioni, nessuna condizione di compatibilit\`a) ha un duale
%	sovradeterminato (al pi\`u una soluzione, una condizione di compatibilit\`a).
%\end{oss}
%
%
%\subsection{Sistema sovradeterminato ($m=3$, $n=r=2$)}
%
%\textit{Primale:}\quad
%$\bm{A}\bm{x}=\bm{b}$\; con\;
%$\bm{A}=\bigl[\begin{smallmatrix}1 & 0\\ 0 & 1\\ 1 & 1\end{smallmatrix}\bigr]$,\;
%$\bm{b}=(1\,\;2\,\;3)^\top$ (compatibile).\quad
%Soluzione unica: $\bm{x}=(1\,\;2)^\top$.
%
%\medskip\noindent
%\textit{Duale:}\quad
%$\bm{A}^\top\bm{y}=\bm{c}$\; con\;
%$\bm{A}^\top=\bigl[\begin{smallmatrix}1 & 0 & 1\\ 0 & 1 & 1\end{smallmatrix}\bigr]$,\;
%$\bm{c}=(c_1\,\;c_2)^\top$.
%
%\noindent Il duale \`e un sistema sottodeterminato ($2$ equazioni, $3$ incognite),
%compatibile per ogni $\bm{c}$. Per $\bm{c}=(3\,\;4)^\top$:
%\begin{equation*}
%	\bm{y}
%	= \underbrace{\begin{Bmatrix} 3 \\ 4 \\ 0 \end{Bmatrix}}_{\bm{y}_0}
%	+ \underbrace{\begin{Bmatrix} -1 \\ -1 \\ 1 \end{Bmatrix}}_{\bm{Y}_\beta} \beta.
%\end{equation*}
%
%\noindent\textit{Verifica dell'identit\`a bilineare} (con $\beta=1$, cio\`e
%$\bm{y}=(2\,\;3\,\;1)^\top$):
%\begin{equation*}
%	\bm{y}^\top\bm{b} = (2)(1)+(3)(2)+(1)(3) = 11, \qquad
%	\bm{x}^\top\bm{c} = (1)(3)+(2)(4) = 11. \quad\checkmark
%\end{equation*}
%
%\noindent\textit{Sottospazi:}
%\begin{equation*}
%	\begin{aligned}
%		&\Nc(\bm{A}) = \{\bm{0}\} \subset \mathbb{R}^2, \quad
%		&&\Ic(\bm{A}) = \Span\{(1\,\;0\,\;1)^\top,\;(0\,\;1\,\;1)^\top\} \subset \mathbb{R}^3, \\
%		&\Nc(\bm{A}^\top) = \Span\{(1\,\;1\,\;-1)^\top\} \subset \mathbb{R}^3, \quad
%		&&\Ic(\bm{A}^\top) = \mathbb{R}^2.
%	\end{aligned}
%\end{equation*}
%Si verifica: $(1\,\;1\,\;-1)\cdot(1\,\;0\,\;1)=0$ e $(1\,\;1\,\;-1)\cdot(0\,\;1\,\;1)=0$, cio\`e
%$\Nc(\bm{A}^\top)\perp\Ic(\bm{A})$, e la somma diretta
%$\mathbb{R}^3=\Nc(\bm{A}^\top)\oplus\Ic(\bm{A})$.
%
%
%\subsection{Sistema degenere ($m=n=3$, $r=2$)}
%
%\textit{Primale:}\quad
%$\bm{A}\bm{x}=\bm{b}$\; con\;
%$\bm{A}=\bigl[\begin{smallmatrix}1 & 0 & 1\\ 0 & 1 & 1\\ 1 & 1 & 2\end{smallmatrix}\bigr]$,\;
%$\bm{b}=(2\,\;1\,\;3)^\top$ (compatibile).\quad
%Soluzione: $\bm{x}=(2\,\;1\,\;0)^\top+(-1\,\;-1\,\;1)^\top\alpha$.
%
%\medskip\noindent
%\textit{Duale:}\quad
%$\bm{A}^\top\bm{y}=\bm{c}$\; con\;
%$\bm{A}^\top=\bigl[\begin{smallmatrix}1 & 0 & 1\\ 0 & 1 & 1\\ 1 & 1 & 2\end{smallmatrix}\bigr]$.\;
%Essendo $\bm{A}$ simmetrica, anche il duale ha la stessa matrice e lo stesso
%rango $r=2$. Il duale \`e anch'esso degenere, con una condizione di
%compatibilit\`a ($c_1+c_2-c_3=0$) e un parametro libero.
%
%Per $\bm{c}=(1\,\;1\,\;2)^\top$ (compatibile):
%\begin{equation*}
%	\bm{y}
%	= \underbrace{\begin{Bmatrix} 1 \\ 1 \\ 0 \end{Bmatrix}}_{\bm{y}_0}
%	+ \underbrace{\begin{Bmatrix} -1 \\ -1 \\ 1 \end{Bmatrix}}_{\bm{Y}_\beta} \beta.
%\end{equation*}
%
%\noindent\textit{Verifica dell'identit\`a bilineare} (con $\alpha=1$, $\beta=0$,
%cio\`e $\bm{x}=(1,\;0,\;1)^\top$, $\bm{y}=(1\,\;1\,\;0)^\top$):
%\begin{equation*}
%	\bm{y}^\top\bm{b} = (1)(2)+(1)(1)+(0)(3) = 3, \qquad
%	\bm{x}^\top\bm{c} = (1)(1)+(0)(1)+(1)(2) = 3. \quad\checkmark
%\end{equation*}
%
%\noindent\textit{Sottospazi:}
%\begin{equation*}
%	\begin{aligned}
%		&\Nc(\bm{A}) = \Span\{(-1\,\;-1\,\;1)^\top\}, \quad
%		&&\Ic(\bm{A}) = \Span\{(1\,\;0\,\;1)^\top,\;(0\,\;1\,\;1)^\top\}, \\
%		&\Nc(\bm{A}^\top) = \Span\{(-1\,\;-1\,\;1)^\top\}, \quad
%		&&\Ic(\bm{A}^\top) = \Span\{(1\,\;0\,\;1)^\top,\;(0\,\;1\,\;1)^\top\}.
%	\end{aligned}
%\end{equation*}
%Essendo $\bm{A}$ simmetrica, i quattro sottospazi formano due sole coppie:
%$\Nc(\bm{A})=\Nc(\bm{A}^\top)$ e $\Ic(\bm{A})=\Ic(\bm{A}^\top)$.
%Primale e duale hanno la stessa struttura.
%
%\begin{oss}
%	In tutti e quattro gli esempi l'identit\`a bilineare
%	$\bm{y}^\top\bm{b}=\bm{x}^\top\bm{c}$ \`e verificata, indipendentemente
%	dalla scelta dei parametri liberi. Ci\`o \`e conseguenza della struttura
%	algebrica e non del caso numerico particolare. Nel testo,
%	questa identit\`a assumer\`a il significato fisico del principio dei lavori
%	virtuali.
%\end{oss}
%
%%-------------------------------------------------------------------------------------------------------------------
%\section{Esercizi proposti}
%%-------------------------------------------------------------------------------------------------------------------
%
%\begin{esercizio} [Formulazione del problema duale]
%	Dato il sistema primale
%	\begin{equation*}
%		\begin{bmatrix} 1 & 3 \\ 2 & 1 \end{bmatrix}
%		\begin{Bmatrix} x_1 \\ x_2 \end{Bmatrix}
%		= \begin{Bmatrix} 5 \\ 4 \end{Bmatrix},
%	\end{equation*}
%	(a) scrivere il corrispondente problema duale
%	$\bm{A}^\top\bm{y}=\bm{c}$ per $\bm{c}=(3\,\;7)^\top$;
%	(b) costruire la rappresentazione schematica della dualit\`a
%	(tabella con $x_j$ in testa, $y_i$ a lato, $b_i$ a destra,
%	$c_j$ in basso); (c) risolvere entrambi i problemi.
%\end{esercizio}
%
%\begin{esercizio} [Combinazioni lineari di colonne e di righe]
%	Per il sistema dell'Esercizio~1, riscrivere il problema primale
%	nella forma $x_1\bm{a}_1+x_2\bm{a}_2=\bm{b}$ e il problema
%	duale nella forma $y_1\bm{a}^1+y_2\bm{a}^2=\bm{c}$, dove
%	$\bm{a}_j$ \`e la $j$-esima colonna e $\bm{a}^i$ la $i$-esima
%	riga di $\bm{A}$. Verificare che le soluzioni trovate
%	nell'Esercizio~1 soddisfano entrambe le rappresentazioni.
%\end{esercizio}
%
%\begin{esercizio} [Verifica dell'identit\`a bilineare]
%	Utilizzando le soluzioni $\bm{x}$ e $\bm{y}$ ottenute
%	nell'Esercizio~1, verificare l'identit\`a bilineare
%	$\bm{y}^\top\bm{b}=\bm{x}^\top\bm{c}$. Ripetere la verifica
%	per il sistema
%	\begin{equation*}
%		\bm{A} = \begin{bmatrix} 1 & 2 \end{bmatrix}, \quad
%		b = 4, \quad \bm{c} = (1\,\;2)^\top,
%	\end{equation*}
%	scegliendo due valori diversi del parametro libero $\alpha$ del
%	problema primale e osservando che l'identit\`a \`e soddisfatta
%	indipendentemente da tale scelta.
%\end{esercizio}
%
%\begin{esercizio} [Significato dei coefficienti]
%	Si consideri il sistema primale
%	\begin{equation*}
%		\begin{bmatrix} 3 & 1 \\ 2 & 4 \end{bmatrix}
%		\begin{Bmatrix} x_1 \\ x_2 \end{Bmatrix}
%		= \begin{Bmatrix} b_1 \\ b_2 \end{Bmatrix}.
%	\end{equation*}
%	(a) Scrivere il problema duale. (b) Considerando il coefficiente
%	$a_{21}=2$: nel problema primale, qual \`e il contributo unitario
%	di $x_1$ all'equazione 2? Nel problema duale, qual \`e il
%	contributo unitario di $y_2$ all'equazione 1? (c) Verificare che
%	lo sviluppo esplicito dell'identit\`a
%	$\bm{y}^\top\bm{A}\bm{x}=\sum_{i,j}y_i\,a_{ij}\,x_j$ mette in
%	evidenza il doppio ruolo di ciascun $a_{ij}$.
%\end{esercizio}
%
%\begin{esercizio} [Sottospazi fondamentali e ortogonalit\`a]
%	Data la matrice
%	\begin{equation*}
%		\bm{A} = \begin{bmatrix} 1 & 0 & 1 \\ 0 & 1 & -1 \end{bmatrix},
%	\end{equation*}
%	(a) determinare $\Nc(\bm{A})$, $\Ic(\bm{A})$,
%	$\Nc(\bm{A}^\top)$, $\Ic(\bm{A}^\top)$ e le rispettive
%	dimensioni; (b) verificare le relazioni di ortogonalit\`a
%	$\Nc(\bm{A})\perp\Ic(\bm{A}^\top)$ e
%	$\Nc(\bm{A}^\top)\perp\Ic(\bm{A})$; (c) verificare le
%	decomposizioni in somma diretta
%	$\mathbb{R}^3=\Nc(\bm{A})\oplus\Ic(\bm{A}^\top)$ e
%	$\mathbb{R}^2=\Nc(\bm{A}^\top)\oplus\Ic(\bm{A})$.
%\end{esercizio}
%
%\begin{esercizio} [Classificazione incrociata primale/duale]
%	Per la matrice dell'Esercizio~5: (a) classificare il sistema
%	primale $\bm{A}\bm{x}=\bm{b}$ e il duale
%	$\bm{A}^\top\bm{y}=\bm{c}$; (b) per
%	$\bm{b}=(3\,\;1)^\top$ risolvere il primale e per
%	$\bm{c}=(1\,\;1\,\;0)^\top$ verificare la compatibilit\`a del
%	duale; (c) verificare l'identit\`a bilineare; (d) commentare la
%	simmetria: il primale sottodeterminato ha un duale
%	sovradeterminato.
%\end{esercizio}
%
%\begin{esercizio} [Sistema degenere simmetrico]
%	Dato il sistema con matrice simmetrica
%	\begin{equation*}
%		\bm{A} = \begin{bmatrix} 1 & 2 & 3 \\ 2 & 4 & 6 \\ 3 & 6 & 9 \end{bmatrix},
%	\end{equation*}
%	(a) determinare il rango e verificare che $\bm{A}$ \`e simmetrica;
%	(b) identificare i quattro sottospazi fondamentali e osservare che
%	$\Nc(\bm{A})=\Nc(\bm{A}^\top)$ e
%	$\Ic(\bm{A})=\Ic(\bm{A}^\top)$;
%	(c) per $\bm{b}=(1,\;2,\;3)^\top$ risolvere il primale, scrivere
%	il duale con $\bm{c}=\bm{b}$ e risolvere anche quest'ultimo;
%	(d) verificare l'identit\`a bilineare e commentare il fatto che
%	primale e duale hanno la stessa struttura.
%	
%	
%	\noindent\textit{Suggerimento per il punto} (d): poich\'e $\bm{A}$
%	\`e simmetrica, primale e duale sono governati dallo stesso
%	operatore; i quattro sottospazi si riducono a due coppie coincidenti
%	e i due problemi condividono la stessa classificazione, le stesse
%	condizioni di compatibilit\`a e lo stesso numero di gradi di
%	libert\`a nella soluzione. Il sistema \`e \textit{autoduale}.
%\end{esercizio}
%	
